\documentclass[11pt]{article}
\usepackage[a4paper,margin=0.72in]{geometry}
\usepackage{amsmath,amssymb,amsthm}
\usepackage{enumitem}
\usepackage{xcolor}
\setlength{\parindent}{0pt}
\setlength{\parskip}{0.55em}
\setlist[enumerate]{leftmargin=*,itemsep=0.25em,topsep=0.25em}
\allowdisplaybreaks

\definecolor{navy}{RGB}{20,35,70}
\definecolor{softblue}{RGB}{235,244,255}
\definecolor{softgreen}{RGB}{238,249,241}
\definecolor{softred}{RGB}{255,241,241}
\definecolor{softyellow}{RGB}{255,249,230}
\definecolor{lineblue}{RGB}{40,95,180}
\definecolor{linegreen}{RGB}{40,130,80}
\definecolor{linered}{RGB}{180,60,60}



\newenvironment{methodbox}[1]{\par\medskip\noindent\textcolor{lineblue}{\textbf{#1}}\par\smallskip}{\par\medskip}
\newenvironment{answerbox}[1]{\par\medskip\noindent\textcolor{linegreen}{\textbf{#1}}\par\smallskip}{\par\medskip}
\newenvironment{warningbox}[1]{\par\medskip\noindent\textcolor{orange!80!black}{\textbf{#1}}\par\smallskip}{\par\medskip}
\newcommand{\R}{\mathbb{R}}
\newcommand{\grad}{\nabla}
\newcommand{\Hess}{\operatorname{Hess}}
\newcommand{\dd}{\,d}
\newcommand{\ve}[1]{\mathbf{#1}}
\newcommand{\crit}{\mathcal{C}}
\newcommand{\pd}{\frac{\partial}}
\newcommand{\norm}[1]{\left\lVert #1\right\rVert}
\newcommand{\ip}[2]{\left\langle #1,#2\right\rangle}

\title{\vspace{-1.5em}\textbf{Multivariable Calculus: Problem Set 7}\\[0.25em]
\large Detailed Assignment-Style Solutions with Diagrams}
\author{Lecturer: Rajendra Bhatia \quad TA: Debabrata Jana}
\date{}

\begin{document}
\maketitle
\vspace{-1.5em}

\begin{methodbox}{How to read these solutions}
Every problem is written in an assignment style: first the recognition cue, then the method, then the calculation, and finally the conclusion.  I do not omit algebraic steps that are needed to justify the answer.  For local extrema in two variables, the standard interior test is based on the Hessian matrix
\[
H_f(x,y)=\begin{pmatrix} f_{xx} & f_{xy} \\ f_{yx} & f_{yy}\end{pmatrix},
\qquad D=f_{xx}f_{yy}-f_{xy}^2.
\]
At an interior critical point, if $D>0$ and $f_{xx}>0$, the point is a local minimum; if $D>0$ and $f_{xx}<0$, the point is a local maximum; if $D<0$, the point is a saddle point; if $D=0$, the test is inconclusive.
\end{methodbox}



\begin{center}
\renewcommand{\arraystretch}{1.55}
\begin{tabular}{c}
\fbox{\textbf{Find $\nabla f=0$}}\\
$\Downarrow$\\
\fbox{\textbf{Compute $D=f_{xx}f_{yy}-f_{xy}^{2}$}}\\
$\swarrow \qquad \Downarrow \qquad \searrow$\\
\fbox{$D>0,\ f_{xx}>0$: local minimum}\quad
\fbox{$D>0,\ f_{xx}<0$: local maximum}\quad
\fbox{$D<0$: saddle point}\\
$\Downarrow$\\
\fbox{$D=0$: the Hessian test is inconclusive; use another method.}
\end{tabular}
\end{center}



\tableofcontents
\newpage

\section{Problem 1: Positive definite matrices}

\subsection*{Problem 1(i)}
Let $A$ be a $2\times 2$ symmetric matrix. Show that $A$ is positive definite if and only if all its diagonal entries and determinant are positive.

\begin{methodbox}{Method}
Write a general symmetric $2\times 2$ matrix as
\[
A=\begin{pmatrix}a&b\\ b&d\end{pmatrix}.
\]
The quadratic form is
\[
q(x,y)=\begin{pmatrix}x&y\end{pmatrix}A\begin{pmatrix}x\\y\end{pmatrix}=ax^2+2bxy+dy^2.
\]
Positive definite means $q(x,y)>0$ for every $(x,y)\neq(0,0)$.
\end{methodbox}

First assume that $A$ is positive definite.  Taking $(x,y)=(1,0)$ gives
\[
q(1,0)=a>0.
\]
Taking $(x,y)=(0,1)$ gives
\[
q(0,1)=d>0.
\]
So both diagonal entries are positive.

It remains to show that $\det A>0$. Since $a>0$, complete the square:
\[
ax^2+2bxy+dy^2
=a\left(x+\frac{b}{a}y\right)^2+\left(d-\frac{b^2}{a}\right)y^2.
\]
Choose $x=-\frac{b}{a}y$ and choose $y\neq0$.  Since $A$ is positive definite,
\[
q\left(-\frac{b}{a}y,y\right)=\left(d-\frac{b^2}{a}\right)y^2>0.
\]
Because $y^2>0$, this implies
\[
d-\frac{b^2}{a}>0.
\]
Multiplying by $a>0$ gives
\[
ad-b^2>0.
\]
Therefore
\[
\det A=ad-b^2>0.
\]

Conversely, assume that
\[
a>0,\qquad d>0,\qquad ad-b^2>0.
\]
We again complete the square:
\[
q(x,y)=a\left(x+\frac{b}{a}y\right)^2+\frac{ad-b^2}{a}y^2.
\]
Here $a>0$ and $\frac{ad-b^2}{a}>0$. Hence both terms are nonnegative. If $q(x,y)=0$, then both terms must be zero. Thus
\[
y=0,
\qquad x+\frac{b}{a}y=0.
\]
Since $y=0$, the second equation gives $x=0$. Therefore $q(x,y)>0$ for every nonzero vector $(x,y)$. Thus $A$ is positive definite.

\begin{answerbox}{Conclusion}
For a symmetric matrix $A=\begin{pmatrix}a&b\\b&d\end{pmatrix}$,
\[
A\text{ is positive definite}\quad\Longleftrightarrow\quad a>0,
\quad d>0,
\quad ad-b^2>0.
\]
\end{answerbox}

\subsection*{Problem 1(ii)}
Let $A$ be a $2\times 2$ symmetric matrix. Show that $A$ is negative definite if and only if all its diagonal entries are negative and determinant is positive.

\begin{methodbox}{Method}
A matrix $A$ is negative definite exactly when $-A$ is positive definite.  So we apply the result from Problem 1(i) to $-A$.
\end{methodbox}

Let
\[
A=\begin{pmatrix}a&b\\b&d\end{pmatrix}.
\]
Then
\[
-A=\begin{pmatrix}-a&-b\\-b&-d\end{pmatrix}.
\]
By Problem 1(i), $-A$ is positive definite if and only if
\[
-a>0,
\qquad -d>0,
\qquad \det(-A)>0.
\]
The first two inequalities mean
\[
a<0,
\qquad d<0.
\]
For a $2\times 2$ matrix,
\[
\det(-A)=(-a)(-d)-(-b)^2=ad-b^2=\det A.
\]
Hence $\det(-A)>0$ is the same as $\det A>0$.

\begin{answerbox}{Conclusion}
A symmetric $2\times 2$ matrix is negative definite if and only if its diagonal entries are negative and its determinant is positive:
\[
a<0,
\qquad d<0,
\qquad ad-b^2>0.
\]
\end{answerbox}

\subsection*{Problem 1(iii)}
Let $A$ be an $n\times n$ positive definite matrix. Show that all its diagonal entries are positive.

\begin{methodbox}{Method}
Use the standard basis vector $e_i$, whose $i$-th coordinate is $1$ and all other coordinates are $0$.
\end{methodbox}

Since $A$ is positive definite,
\[
x^T A x>0
\]
for every nonzero $x\in\R^n$.  Choose
\[
x=e_i=(0,\ldots,0,1,0,\ldots,0)^T.
\]
Then $e_i\neq0$, so
\[
e_i^T A e_i>0.
\]
But $e_i^T A e_i$ selects exactly the $(i,i)$ entry of $A$. Therefore
\[
e_i^T A e_i=a_{ii}.
\]
Thus
\[
a_{ii}>0
\]
for every $i=1,\ldots,n$.

\begin{answerbox}{Conclusion}
Every diagonal entry of a positive definite matrix is positive:
\[
a_{11}>0,
\ a_{22}>0,
\ldots,
\ a_{nn}>0.
\]
\end{answerbox}

\newpage
\section{Problem 2: Local extrema and saddle points}

\begin{methodbox}{General method for Problem 2}
For each differentiable function $f(x,y)$:
\begin{enumerate}
\item Compute the first partial derivatives $f_x$ and $f_y$.
\item Solve the critical point equations $f_x=0$ and $f_y=0$.
\item Compute
\[
D=f_{xx}f_{yy}-f_{xy}^2.
\]
\item Classify each critical point using the Hessian test.
\item If a domain boundary is explicitly given, separate interior critical point classification from boundary behavior.
\end{enumerate}
\end{methodbox}

\subsection*{Problem 2(i)}
\[
f(x,y)=x^2+xy+y^2+y.
\]
Compute the first partial derivatives:
\[
f_x=2x+y,
\qquad
f_y=x+2y+1.
\]
Critical points satisfy
\[
2x+y=0,
\qquad x+2y+1=0.
\]
From the first equation,
\[
y=-2x.
\]
Substitute this into the second equation:
\[
x+2(-2x)+1=0,
\]
so
\[
x-4x+1=0,
\qquad -3x+1=0,
\qquad x=\frac13.
\]
Then
\[
y=-2\left(\frac13\right)=-\frac23.
\]
So the only critical point is
\[
\left(\frac13,-\frac23\right).
\]

Now compute second partial derivatives:
\[
f_{xx}=2,
\qquad f_{yy}=2,
\qquad f_{xy}=1.
\]
Therefore
\[
D=(2)(2)-1^2=4-1=3>0.
\]
Since $D>0$ and $f_{xx}=2>0$, the point is a local minimum.

The value is
\[
f\left(\frac13,-\frac23\right)
=\frac19-\frac29+\frac49-\frac23.
\]
Put everything over $9$:
\[
\frac19-\frac29+\frac49-\frac69
=\frac{1-2+4-6}{9}
=-\frac13.
\]

\begin{answerbox}{Answer}
The function has a local minimum at
\[
\left(\frac13,-\frac23\right)
\]
with value
\[
-\frac13.
\]
There is no local maximum and no saddle point.
\end{answerbox}

\subsection*{Problem 2(ii)}
\[
f(x,y)=x^3+y^2-3x.
\]
Compute the first partial derivatives:
\[
f_x=3x^2-3,
\qquad f_y=2y.
\]
Set them equal to zero:
\[
3x^2-3=0,
\qquad 2y=0.
\]
Thus
\[
x^2=1,
\qquad y=0.
\]
So the critical points are
\[
(1,0),
\qquad (-1,0).
\]

Second partial derivatives are
\[
f_{xx}=6x,
\qquad f_{yy}=2,
\qquad f_{xy}=0.
\]
Hence
\[
D=(6x)(2)-0=12x.
\]
At $(1,0)$,
\[
D=12>0,
\qquad f_{xx}=6>0.
\]
So $(1,0)$ is a local minimum.  Its value is
\[
f(1,0)=1^3+0^2-3(1)=-2.
\]
At $(-1,0)$,
\[
D=-12<0.
\]
So $(-1,0)$ is a saddle point.  Its value is
\[
f(-1,0)=(-1)^3+0^2-3(-1)=-1+3=2.
\]

\begin{answerbox}{Answer}
Local minimum: $(1,0)$ with value $-2$.  Saddle point: $(-1,0)$ with value $2$.  There is no local maximum.
\end{answerbox}

\subsection*{Problem 2(iii)}
\[
f(x,y)=(x-y)(1-xy).
\]
First expand the function because derivatives are easier after expansion:
\[
f(x,y)=x-y-x^2y+xy^2.
\]
Then
\[
f_x=1-2xy+y^2,
\qquad
f_y=-1-x^2+2xy.
\]
Critical points satisfy
\[
y^2-2xy+1=0,
\qquad
2xy-x^2-1=0.
\]
From the second equation,
\[
2xy=x^2+1.
\]
Substitute this into the first equation:
\[
y^2-(x^2+1)+1=0.
\]
Therefore
\[
y^2-x^2=0.
\]
So
\[
y=x
\qquad\text{or}\qquad y=-x.
\]
If $y=x$, then the equation $2xy=x^2+1$ becomes
\[
2x^2=x^2+1,
\]
so
\[
x^2=1.
\]
Thus
\[
(x,y)=(1,1),\qquad (-1,-1).
\]
If $y=-x$, then
\[
2x(-x)=x^2+1,
\]
so
\[
-2x^2=x^2+1,
\qquad -3x^2=1,
\]
which has no real solution. Hence the only critical points are
\[
(1,1),
\qquad (-1,-1).
\]

Now compute the Hessian entries:
\[
f_{xx}=-2y,
\qquad f_{yy}=2x,
\qquad f_{xy}=-2x+2y.
\]
At $(1,1)$,
\[
f_{xx}=-2,
\qquad f_{yy}=2,
\qquad f_{xy}=0.
\]
Thus
\[
D=(-2)(2)-0=-4<0.
\]
Therefore $(1,1)$ is a saddle point.

At $(-1,-1)$,
\[
f_{xx}=2,
\qquad f_{yy}=-2,
\qquad f_{xy}=0.
\]
Thus
\[
D=(2)(-2)-0=-4<0.
\]
Therefore $(-1,-1)$ is also a saddle point.

\begin{answerbox}{Answer}
Both critical points are saddle points:
\[
(1,1),
\qquad (-1,-1).
\]
There are no local maxima and no local minima.
\end{answerbox}

\subsection*{Problem 2(iv)}
\[
f(x,y)=y^3+3x^2y-6x^2-6y^2+2.
\]
Compute first partial derivatives:
\[
f_x=6xy-12x=6x(y-2),
\]
\[
f_y=3y^2+3x^2-12y=3(x^2+y^2-4y).
\]
Critical points satisfy
\[
6x(y-2)=0,
\qquad 3(x^2+y^2-4y)=0.
\]
The first equation gives
\[
x=0
\qquad\text{or}\qquad y=2.
\]
The second equation gives
\[
x^2+y^2-4y=0.
\]

Case 1: $x=0$.  Then
\[
y^2-4y=0,
\]
so
\[
y(y-4)=0.
\]
Thus
\[
y=0\quad\text{or}\quad y=4.
\]
This gives
\[
(0,0),
\qquad (0,4).
\]

Case 2: $y=2$.  Then
\[
x^2+2^2-4(2)=0,
\]
so
\[
x^2+4-8=0,
\qquad x^2=4.
\]
Thus
\[
x=2\quad\text{or}\quad x=-2.
\]
This gives
\[
(2,2),
\qquad (-2,2).
\]

Second partial derivatives are
\[
f_{xx}=6y-12,
\qquad f_{yy}=6y-12,
\qquad f_{xy}=6x.
\]
At $(0,0)$,
\[
f_{xx}=-12,
\quad f_{yy}=-12,
\quad f_{xy}=0.
\]
So
\[
D=(-12)(-12)-0=144>0,
\qquad f_{xx}<0.
\]
Thus $(0,0)$ is a local maximum. Its value is
\[
f(0,0)=2.
\]
At $(0,4)$,
\[
f_{xx}=12,
\quad f_{yy}=12,
\quad f_{xy}=0.
\]
So
\[
D=144>0,
\qquad f_{xx}>0.
\]
Thus $(0,4)$ is a local minimum. Its value is
\[
f(0,4)=4^3-6(4^2)+2=64-96+2=-30.
\]
At $(2,2)$,
\[
f_{xx}=0,
\quad f_{yy}=0,
\quad f_{xy}=12.
\]
Thus
\[
D=0\cdot0-12^2=-144<0.
\]
So $(2,2)$ is a saddle point.

At $(-2,2)$,
\[
f_{xx}=0,
\quad f_{yy}=0,
\quad f_{xy}=-12.
\]
Thus
\[
D=0\cdot0-(-12)^2=-144<0.
\]
So $(-2,2)$ is a saddle point.

\begin{answerbox}{Answer}
Local maximum: $(0,0)$ with value $2$.  Local minimum: $(0,4)$ with value $-30$.  Saddle points: $(2,2)$ and $(-2,2)$.
\end{answerbox}

\subsection*{Problem 2(v)}
\[
f(x,y)=xy(1-x-y).
\]
Expand:
\[
f(x,y)=xy-x^2y-xy^2.
\]
Compute first partial derivatives:
\[
f_x=y-2xy-y^2=y(1-2x-y),
\]
\[
f_y=x-x^2-2xy=x(1-x-2y).
\]
Critical points satisfy
\[
y(1-2x-y)=0,
\qquad x(1-x-2y)=0.
\]
We solve by cases.

If $x=0$ and $y=0$, we get $(0,0)$.

If $x=0$ and $1-2x-y=0$, then
\[
1-y=0,
\]
so $y=1$, giving $(0,1)$.

If $y=0$ and $1-x-2y=0$, then
\[
1-x=0,
\]
so $x=1$, giving $(1,0)$.

If
\[
1-2x-y=0,
\qquad 1-x-2y=0,
\]
then
\[
2x+y=1,
\qquad x+2y=1.
\]
Subtract the second equation from the first:
\[
x-y=0,
\]
so $x=y$.  Substituting into $2x+y=1$ gives
\[
3x=1,
\qquad x=\frac13.
\]
Thus
\[
y=\frac13.
\]
So the fourth critical point is
\[
\left(\frac13,\frac13\right).
\]

Second partial derivatives are
\[
f_{xx}=-2y,
\qquad f_{yy}=-2x,
\qquad f_{xy}=1-2x-2y.
\]
At $(0,0)$,
\[
H=\begin{pmatrix}0&1\\1&0\end{pmatrix},
\qquad D=-1<0.
\]
So $(0,0)$ is a saddle point.

At $(1,0)$,
\[
f_{xx}=0,
\quad f_{yy}=-2,
\quad f_{xy}=-1.
\]
Thus
\[
D=0(-2)-(-1)^2=-1<0.
\]
So $(1,0)$ is a saddle point.

At $(0,1)$,
\[
f_{xx}=-2,
\quad f_{yy}=0,
\quad f_{xy}=-1.
\]
Thus
\[
D=(-2)(0)-(-1)^2=-1<0.
\]
So $(0,1)$ is a saddle point.

At $\left(\frac13,\frac13\right)$,
\[
f_{xx}=-\frac23,
\quad f_{yy}=-\frac23,
\quad f_{xy}=1-\frac23-\frac23=-\frac13.
\]
Therefore
\[
D=\left(-\frac23\right)\left(-\frac23\right)-\left(-\frac13\right)^2
=\frac49-\frac19
=\frac13>0.
\]
Since $D>0$ and $f_{xx}<0$, this point is a local maximum.

Its value is
\[
f\left(\frac13,\frac13\right)
=\frac13\cdot\frac13\left(1-\frac13-\frac13\right)
=\frac19\cdot\frac13
=\frac1{27}.
\]

\begin{answerbox}{Answer}
Local maximum: $\left(\frac13,\frac13\right)$ with value $\frac1{27}$.  Saddle points: $(0,0)$, $(1,0)$, and $(0,1)$.  There is no local minimum.
\end{answerbox}

\subsection*{Problem 2(vi)}
\[
f(x,y)=e^x\cos y.
\]
Compute first partial derivatives:
\[
f_x=e^x\cos y,
\qquad f_y=-e^x\sin y.
\]
A critical point would require
\[
e^x\cos y=0,
\qquad -e^x\sin y=0.
\]
Since $e^x>0$ for every real $x$, these equations reduce to
\[
\cos y=0,
\qquad \sin y=0.
\]
But no real number $y$ satisfies both $\sin y=0$ and $\cos y=0$ at the same time.

\begin{answerbox}{Answer}
There are no critical points. Therefore the function has no local maxima, no local minima, and no saddle points in the usual interior critical-point sense.
\end{answerbox}

\subsection*{Problem 2(vii)}
\[
f(x,y)=(x^2+y^2)e^{y^2-x^2}.
\]
Let
\[
r^2=x^2+y^2.
\]
Then
\[
f=r^2e^{y^2-x^2}.
\]
Using the product rule,
\[
f_x=2x e^{y^2-x^2}+r^2e^{y^2-x^2}(-2x).
\]
Factor out $2xe^{y^2-x^2}$:
\[
f_x=2xe^{y^2-x^2}(1-r^2).
\]
Similarly,
\[
f_y=2ye^{y^2-x^2}+r^2e^{y^2-x^2}(2y)
=2ye^{y^2-x^2}(1+r^2).
\]
Critical points satisfy
\[
2xe^{y^2-x^2}(1-r^2)=0,
\qquad
2ye^{y^2-x^2}(1+r^2)=0.
\]
Since $e^{y^2-x^2}>0$ and $1+r^2>0$, the second equation gives
\[
y=0.
\]
With $y=0$, we have $r^2=x^2$, and the first equation becomes
\[
2xe^{-x^2}(1-x^2)=0.
\]
Thus
\[
x=0
\qquad\text{or}\qquad x^2=1.
\]
Therefore the critical points are
\[
(0,0),
\qquad (1,0),
\qquad (-1,0).
\]

At $(0,0)$, since
\[
f(x,y)=(x^2+y^2)e^{y^2-x^2}\ge 0
\]
and
\[
f(0,0)=0,
\]
we immediately get a local minimum at $(0,0)$.  In fact it is also a global minimum.

At $(1,0)$ and $(-1,0)$, we use the Hessian. From
\[
f_x=2xe^{y^2-x^2}(1-r^2),
\qquad
f_y=2ye^{y^2-x^2}(1+r^2),
\]
one obtains at $(\pm1,0)$:
\[
f_{xx}=-\frac4e,
\qquad f_{yy}=\frac4e,
\qquad f_{xy}=0.
\]
Therefore
\[
D=\left(-\frac4e\right)\left(\frac4e\right)-0=-\frac{16}{e^2}<0.
\]
So both $(1,0)$ and $(-1,0)$ are saddle points.

\begin{answerbox}{Answer}
Local minimum: $(0,0)$ with value $0$.  Saddle points: $(1,0)$ and $(-1,0)$.  There is no local maximum.
\end{answerbox}

\subsection*{Problem 2(viii)}
\[
f(x,y)=e^y(y^2-x^2).
\]
Compute first partial derivatives:
\[
f_x=e^y(-2x)=-2xe^y,
\]
\[
f_y=e^y(y^2-x^2)+e^y(2y)=e^y(y^2+2y-x^2).
\]
Critical points satisfy
\[
-2xe^y=0,
\qquad e^y(y^2+2y-x^2)=0.
\]
Since $e^y>0$, the first equation gives
\[
x=0.
\]
Substitute $x=0$ into the second equation:
\[
y^2+2y=0.
\]
Thus
\[
y(y+2)=0,
\]
so
\[
y=0
\qquad\text{or}\qquad y=-2.
\]
The critical points are
\[
(0,0),
\qquad (0,-2).
\]

Second partial derivatives are
\[
f_{xx}=-2e^y,
\qquad f_{xy}=-2xe^y,
\]
\[
f_{yy}=e^y(y^2+2y-x^2)+e^y(2y+2)
=e^y(y^2+4y+2-x^2).
\]
At $(0,0)$,
\[
f_{xx}=-2,
\quad f_{yy}=2,
\quad f_{xy}=0.
\]
Therefore
\[
D=(-2)(2)-0=-4<0.
\]
So $(0,0)$ is a saddle point.

At $(0,-2)$,
\[
f_{xx}=-2e^{-2},
\quad f_{xy}=0,
\quad f_{yy}=e^{-2}(4-8+2)=-2e^{-2}.
\]
Thus
\[
D=(-2e^{-2})(-2e^{-2})-0=4e^{-4}>0.
\]
Since $f_{xx}<0$, $(0,-2)$ is a local maximum.
Its value is
\[
f(0,-2)=e^{-2}((-2)^2-0)=\frac4{e^2}.
\]

\begin{answerbox}{Answer}
Local maximum: $(0,-2)$ with value $4/e^2$.  Saddle point: $(0,0)$.  There is no local minimum.
\end{answerbox}

\subsection*{Problem 2(ix)}
\[
f(x,y)=y^2-2y\cos x,
\qquad 1\le x\le 7.
\]
Complete the square in $y$:
\[
f(x,y)=\left(y-\cos x\right)^2-\cos^2x.
\]
This form is useful because for each fixed $x$, the smallest value in the $y$-direction occurs at $y=\cos x$.

For interior critical points, compute partial derivatives:
\[
f_x=2y\sin x,
\qquad f_y=2y-2\cos x.
\]
Set them equal to zero:
\[
2y\sin x=0,
\qquad 2y-2\cos x=0.
\]
The second equation gives
\[
y=\cos x.
\]
Substitute into the first equation:
\[
2\cos x\sin x=0.
\]
Thus
\[
\cos x=0
\qquad\text{or}\qquad \sin x=0.
\]
For $1<x<7$, the relevant values are
\[
x=\frac\pi2,
\quad \pi,
\quad \frac{3\pi}{2},
\quad 2\pi.
\]
Hence the interior critical points are
\[
\left(\frac\pi2,0\right),
\quad (\pi,-1),
\quad \left(\frac{3\pi}{2},0\right),
\quad (2\pi,1).
\]

Second partial derivatives are
\[
f_{xx}=2y\cos x,
\qquad f_{xy}=2\sin x,
\qquad f_{yy}=2.
\]
At $(\pi,-1)$,
\[
f_{xx}=2(-1)(-1)=2,
\qquad f_{xy}=0,
\qquad f_{yy}=2.
\]
So
\[
D=2\cdot2-0=4>0,
\qquad f_{xx}>0.
\]
Thus $(\pi,-1)$ is a local minimum, with value
\[
f(\pi,-1)=(-1)^2-2(-1)\cos\pi=1+2(-1)=-1.
\]
At $(2\pi,1)$,
\[
f_{xx}=2(1)(1)=2,
\qquad f_{xy}=0,
\qquad f_{yy}=2.
\]
So again it is a local minimum, with value
\[
f(2\pi,1)=1-2(1)(1)=-1.
\]
At $\left(\frac\pi2,0\right)$,
\[
f_{xx}=0,
\qquad f_{xy}=2,
\qquad f_{yy}=2.
\]
Thus
\[
D=0\cdot2-2^2=-4<0.
\]
So $\left(\frac\pi2,0\right)$ is a saddle point.

At $\left(\frac{3\pi}{2},0\right)$,
\[
f_{xx}=0,
\qquad f_{xy}=-2,
\qquad f_{yy}=2.
\]
Thus
\[
D=0\cdot2-(-2)^2=-4<0.
\]
So $\left(\frac{3\pi}{2},0\right)$ is a saddle point.

\begin{warningbox}{Boundary note}
If local extrema are interpreted relative to the closed strip $[1,7]\times\R$, then $x=1$ is a boundary.  The point $(1,\cos1)$ is also a one-sided local minimum because $x\mapsto-\cos^2x$ is increasing at $x=1$.  The right boundary $x=7$ gives no local extremum at $(7,\cos7)$ because moving left along the minimizing curve decreases the value while moving in the $y$-direction increases the value.
\end{warningbox}

\begin{answerbox}{Interior answer}
Interior local minima: $(\pi,-1)$ and $(2\pi,1)$, both with value $-1$. Interior saddle points: $\left(\frac\pi2,0\right)$ and $\left(\frac{3\pi}{2},0\right)$. There is no interior local maximum.
\end{answerbox}

\subsection*{Problem 2(x)}
\[
f(x,y)=\sin x\sin y,
\qquad -\pi\le x\le\pi,
\qquad -\pi\le y\le\pi.
\]
First classify the interior critical points.  Compute
\[
f_x=\cos x\sin y,
\qquad f_y=\sin x\cos y.
\]
Set them equal to zero:
\[
\cos x\sin y=0,
\qquad \sin x\cos y=0.
\]
In the interior $(-\pi,\pi)^2$, the solutions are:
\[
(0,0),
\]
and
\[
\left(\frac\pi2,\frac\pi2\right),
\quad
\left(-\frac\pi2,-\frac\pi2\right),
\quad
\left(\frac\pi2,-\frac\pi2\right),
\quad
\left(-\frac\pi2,\frac\pi2\right).
\]
Second partial derivatives are
\[
f_{xx}=-\sin x\sin y,
\qquad f_{yy}=-\sin x\sin y,
\qquad f_{xy}=\cos x\cos y.
\]
At $(0,0)$,
\[
f_{xx}=0,
\quad f_{yy}=0,
\quad f_{xy}=1,
\]
so
\[
D=0\cdot0-1^2=-1<0.
\]
Thus $(0,0)$ is a saddle point.

At $\left(\frac\pi2,\frac\pi2\right)$ and $\left(-\frac\pi2,-\frac\pi2\right)$,
\[
\sin x\sin y=1.
\]
Therefore
\[
f_{xx}=-1,
\quad f_{yy}=-1,
\quad f_{xy}=0,
\]
so
\[
D=1>0,
\qquad f_{xx}<0.
\]
These are local maxima, both with value $1$.

At $\left(\frac\pi2,-\frac\pi2\right)$ and $\left(-\frac\pi2,\frac\pi2\right)$,
\[
\sin x\sin y=-1.
\]
Therefore
\[
f_{xx}=1,
\quad f_{yy}=1,
\quad f_{xy}=0,
\]
so
\[
D=1>0,
\qquad f_{xx}>0.
\]
These are local minima, both with value $-1$.

\begin{warningbox}{Boundary note for the closed square}
On the closed square, weak one-sided boundary extrema also occur because $\sin(\pm\pi)=0$.  For example, on the edge $x=\pi$, points with $0<y<\pi$ are weak local minima and points with $-\pi<y<0$ are weak local maxima. Similar one-sided statements hold on the other edges.  The main interior critical-point classification is the one above.
\end{warningbox}

\begin{answerbox}{Interior answer}
Local maxima: $\left(\frac\pi2,\frac\pi2\right)$ and $\left(-\frac\pi2,-\frac\pi2\right)$, both with value $1$.  Local minima: $\left(\frac\pi2,-\frac\pi2\right)$ and $\left(-\frac\pi2,\frac\pi2\right)$, both with value $-1$.  Saddle point: $(0,0)$.
\end{answerbox}

\newpage
\section{Problem 3: Extreme values and least squares}

\subsection*{Problem 3(i)}
Find all local and global extreme values for
\[
f(x,y)=xy(1-x^2-y^2)
\]
on
\[
[0,1]^2.
\]

\begin{methodbox}{Method}
Because the domain is closed and bounded, a continuous function has global maximum and minimum values by the extreme value theorem. We check: interior critical points and boundary pieces.
\end{methodbox}

Compute first partial derivatives:
\[
f_x=y(1-x^2-y^2)+xy(-2x)
=y(1-3x^2-y^2),
\]
\[
f_y=x(1-x^2-y^2)+xy(-2y)
=x(1-x^2-3y^2).
\]
In the interior $(0,1)^2$, we have $x>0$ and $y>0$. Therefore the critical point equations reduce to
\[
1-3x^2-y^2=0,
\qquad
1-x^2-3y^2=0.
\]
Equivalently,
\[
3x^2+y^2=1,
\qquad x^2+3y^2=1.
\]
Subtract the second equation from the first:
\[
2x^2-2y^2=0.
\]
Thus
\[
x^2=y^2.
\]
Since $x,y>0$ in the interior, this gives
\[
x=y.
\]
Substitute $y=x$ into $3x^2+y^2=1$:
\[
3x^2+x^2=1,
\qquad 4x^2=1,
\qquad x=\frac12.
\]
Thus
\[
y=\frac12.
\]
The only interior critical point is
\[
\left(\frac12,\frac12\right).
\]
Its value is
\[
f\left(\frac12,\frac12\right)
=\frac12\cdot\frac12\left(1-\frac14-\frac14\right)
=\frac14\cdot\frac12
=\frac18.
\]

Now examine the boundary.

On $x=0$:
\[
f(0,y)=0.
\]
On $y=0$:
\[
f(x,0)=0.
\]
On $x=1$:
\[
f(1,y)=y(1-1-y^2)=-y^3,
\qquad 0\le y\le1.
\]
So on this edge the minimum is $-1$ at $y=1$, and the maximum is $0$ at $y=0$.

On $y=1$:
\[
f(x,1)=x(1-x^2-1)=-x^3,
\qquad 0\le x\le1.
\]
So on this edge the minimum is $-1$ at $x=1$, and the maximum is $0$ at $x=0$.

Therefore the global minimum is
\[
f(1,1)=-1.
\]
To prove that the interior value $\frac18$ is the global maximum, observe that if $x^2+y^2>1$, then
\[
1-x^2-y^2<0,
\]
so
\[
f(x,y)\le0.
\]
Thus positive maxima can only occur where $x^2+y^2\le1$.  Let
\[
s=x^2+y^2.
\]
By $2xy\le x^2+y^2=s$, we get
\[
xy\le\frac{s}{2}.
\]
Therefore
\[
f(x,y)=xy(1-s)\le\frac{s}{2}(1-s).
\]
For $0\le s\le1$,
\[
\frac{s}{2}(1-s)
\]
is maximized at $s=\frac12$, where its value is
\[
\frac12\cdot\frac12\cdot\frac12=\frac18.
\]
Equality requires
\[
x=y,
\qquad x^2+y^2=\frac12,
\]
which gives
\[
x=y=\frac12.
\]



\begin{center}
\renewcommand{\arraystretch}{1.4}
\begin{tabular}{|c|c|c|}
\hline
\textbf{Region in $[0,1]^2$} & \textbf{Sign of $1-x^2-y^2$} & \textbf{Sign of $f=xy(1-x^2-y^2)$}\\
\hline
Inside the quarter disk $x^2+y^2<1$ & Positive & Positive, except on axes where $f=0$\\
\hline
On the curve $x^2+y^2=1$ & Zero & Zero\\
\hline
Outside the quarter disk $x^2+y^2>1$ & Negative & Negative, since $x,y\ge0$\\
\hline
Special points & $\left(\frac12,\frac12\right)$ gives max $\frac18$ & $(1,1)$ gives min $-1$\\
\hline
\end{tabular}
\end{center}



\begin{answerbox}{Answer}
Global maximum:
\[
\frac18\quad\text{at}\quad \left(\frac12,\frac12\right).
\]
Global minimum:
\[
-1\quad\text{at}\quad (1,1).
\]
There are weak boundary local minima on the axes away from $(1,0)$ and $(0,1)$, because $f=0$ there and nearby interior values in the quarter disk are nonnegative. The point $\left(\frac12,\frac12\right)$ is the only interior local maximum.
\end{answerbox}

\subsection*{Problem 3(ii)}
Let
\[
f(x,y)=3x^4-4x^2y+y^2.
\]
Show that on every line $y=mx$, the function has a minimum at $(0,0)$, but $(0,0)$ is not a point of local minimum in $\R^2$.

Along the line
\[
y=mx,
\]
we substitute $y=mx$ into $f$:
\[
f(x,mx)=3x^4-4x^2(mx)+(mx)^2.
\]
Therefore
\[
f(x,mx)=3x^4-4mx^3+m^2x^2.
\]
Factor out $x^2$:
\[
f(x,mx)=x^2(3x^2-4mx+m^2).
\]
At $x=0$, this is $0$.  For fixed $m\ne0$, the leading term near $x=0$ is
\[
m^2x^2,
\]
which is positive for $x\ne0$ sufficiently close to $0$.  Therefore $x=0$ is a local minimum of the restricted one-variable function $x\mapsto f(x,mx)$.

If $m=0$, then $y=0$, and
\[
f(x,0)=3x^4\ge0,
\]
with equality at $x=0$.  Hence $x=0$ is also a minimum along the line $y=0$.

Now show that $(0,0)$ is not a local minimum in the full plane.  Complete the square:
\[
f(x,y)=y^2-4x^2y+3x^4.
\]
Notice that
\[
y^2-4x^2y+4x^4=(y-2x^2)^2.
\]
Thus
\[
f(x,y)=(y-2x^2)^2-x^4.
\]
Now approach $(0,0)$ along the curve
\[
y=2x^2.
\]
Then
\[
f(x,2x^2)=(2x^2-2x^2)^2-x^4=-x^4.
\]
For every $x\ne0$ close to $0$,
\[
-x^4<0=f(0,0).
\]
So arbitrarily close to $(0,0)$, the function is smaller than its value at $(0,0)$.

\begin{answerbox}{Answer}
On every straight line $y=mx$, the origin is a local minimum of the restricted one-variable function.  But in the full plane, along the curve $y=2x^2$, the values are $-x^4<0$. Hence $(0,0)$ is not a local minimum in $\R^2$.
\end{answerbox}

\subsection*{Problem 3(iii)}
Given distinct real numbers $a_1,\ldots,a_n$ and real numbers $b_1,\ldots,b_n$, find the straight line that minimizes
\[
f(x,y)=\sum_{i=1}^n (x+ya_i-b_i)^2.
\]

\begin{methodbox}{Interpretation}
The line has the form
\[
b=x+ya.
\]
Here $x$ is the intercept and $y$ is the slope. The expression $x+ya_i-b_i$ is the vertical error at the data point $(a_i,b_i)$.
\end{methodbox}

Define the following sums:
\[
S_a=\sum_{i=1}^n a_i,
\qquad
S_b=\sum_{i=1}^n b_i,
\qquad
S_{aa}=\sum_{i=1}^n a_i^2,
\qquad
S_{ab}=\sum_{i=1}^n a_i b_i.
\]
Compute the partial derivatives.

First,
\[
f_x=2\sum_{i=1}^n (x+ya_i-b_i).
\]
Set $f_x=0$:
\[
\sum_{i=1}^n (x+ya_i-b_i)=0.
\]
Therefore
\[
nx+y\sum_{i=1}^n a_i-\sum_{i=1}^n b_i=0,
\]
so
\[
nx+yS_a=S_b.
\]

Second,
\[
f_y=2\sum_{i=1}^n a_i(x+ya_i-b_i).
\]
Set $f_y=0$:
\[
\sum_{i=1}^n a_i(x+ya_i-b_i)=0.
\]
Therefore
\[
x\sum_{i=1}^n a_i+y\sum_{i=1}^n a_i^2-\sum_{i=1}^n a_i b_i=0,
\]
so
\[
xS_a+yS_{aa}=S_{ab}.
\]
Thus $(x,y)$ must solve the normal equations
\[
\begin{pmatrix}
n&S_a\\
S_a&S_{aa}
\end{pmatrix}
\begin{pmatrix}x\\y\end{pmatrix}
=
\begin{pmatrix}S_b\\S_{ab}\end{pmatrix}.
\]
The determinant is
\[
\Delta=nS_{aa}-S_a^2.
\]
Since the $a_i$ are not all equal, and in fact are distinct,
\[
\Delta=n\sum_{i=1}^n (a_i-\bar a)^2>0,
\]
where
\[
\bar a=\frac1n\sum_{i=1}^n a_i.
\]
Therefore the solution is unique. By solving the $2\times2$ linear system,
\[
x=\frac{S_bS_{aa}-S_aS_{ab}}{nS_{aa}-S_a^2},
\]
\[
y=\frac{nS_{ab}-S_aS_b}{nS_{aa}-S_a^2}.
\]

\begin{answerbox}{Answer}
The least-squares line is
\[
b=x+ya,
\]
where
\[
x=\frac{S_bS_{aa}-S_aS_{ab}}{nS_{aa}-S_a^2},
\qquad
 y=\frac{nS_{ab}-S_aS_b}{nS_{aa}-S_a^2}.
\]
\end{answerbox}

\newpage
\section{Problem 4: Detailed study of $xy(1-x^2-y^2)$ on the square}

Let
\[
f(x,y)=xy(1-x^2-y^2)
\]
on the square
\[
[0,1]^2.
\]
This is the same function as in Problem 3(i), but now the question asks for the classification more systematically.

\subsection*{Problem 4(i): Critical points}
We already computed
\[
f_x=y(1-3x^2-y^2),
\qquad
f_y=x(1-x^2-3y^2).
\]
In the open square $(0,1)^2$, both $x$ and $y$ are positive, so the critical equations are
\[
1-3x^2-y^2=0,
\qquad
1-x^2-3y^2=0.
\]
Solving these gives
\[
x=y=\frac12.
\]
Therefore the only interior critical point is
\[
\left(\frac12,\frac12\right).
\]
If one solves $\grad f=0$ in the whole plane first, the full critical set in $\R^2$ is
\[
(0,0),\ (1,0),\ (-1,0),\ (0,1),\ (0,-1),
\]
\[
\left(\frac12,\frac12\right),
\left(\frac12,-\frac12\right),
\left(-\frac12,\frac12\right),
\left(-\frac12,-\frac12\right).
\]
Among these, the points lying in $[0,1]^2$ are
\[
(0,0),
\quad (1,0),
\quad (0,1),
\quad \left(\frac12,\frac12\right).
\]
But for the ordinary interior Hessian test on the square, only $\left(\frac12,\frac12\right)$ is an interior critical point.

\subsection*{Problem 4(ii): Local classification on $(0,1)^2$}
Compute the second partial derivatives:
\[
f_{xx}=-6xy,
\qquad
f_{yy}=-6xy,
\qquad
f_{xy}=1-3x^2-3y^2.
\]
At $\left(\frac12,\frac12\right)$,
\[
f_{xx}=-6\cdot\frac12\cdot\frac12=-\frac32,
\]
\[
f_{yy}=-\frac32,
\]
\[
f_{xy}=1-3\cdot\frac14-3\cdot\frac14
=1-\frac34-\frac34
=-\frac12.
\]
Therefore
\[
D=\left(-\frac32\right)\left(-\frac32\right)-\left(-\frac12\right)^2
=\frac94-\frac14
=2>0.
\]
Since $f_{xx}<0$, the point is a local maximum.

\begin{answerbox}{Answer on $(0,1)^2$}
The only interior critical point is
\[
\left(\frac12,\frac12\right),
\]
and it is a local maximum with value
\[
\frac18.
\]
There are no interior local minima and no interior saddle points.
\end{answerbox}

\subsection*{Problem 4(iii): Local behavior on $[0,1]^2$}
On the closed square, boundary points must be understood using relative neighborhoods inside the square.

The sign of $f$ is controlled by
\[
f(x,y)=xy(1-x^2-y^2).
\]
Since $x\ge0$ and $y\ge0$ on the square, the sign depends on
\[
1-x^2-y^2.
\]
Thus:
\[
f>0\quad\text{inside the quarter disk }x^2+y^2<1,
\]
\[
f=0\quad\text{on the axes and on }x^2+y^2=1,
\]
\[
f<0\quad\text{outside the quarter disk, within the square.}
\]

Therefore:
\begin{enumerate}
\item The interior point $\left(\frac12,\frac12\right)$ is a local maximum.
\item Points on the axes $\{(x,0):0\le x<1\}$ and $\{(0,y):0\le y<1\}$ are weak local minima, because $f=0$ there and nearby points sufficiently close to them inside the square have $f\ge0$.
\item The point $(1,1)$ is a local minimum and also the global minimum.
\item Points on the quarter circle $x^2+y^2=1$ with $x>0,y>0$ are not local extrema: nearby points inside the circle give $f>0$, while nearby points outside the circle give $f<0$.
\item The points $(1,0)$ and $(0,1)$ are also not local extrema because values of both signs occur arbitrarily close to them inside the square.
\end{enumerate}

\begin{answerbox}{Answer on $[0,1]^2$}
Local maximum: $\left(\frac12,\frac12\right)$.  Weak boundary local minima: the axes segments $\{(x,0):0\le x<1\}$ and $\{(0,y):0\le y<1\}$, together with the corner $(1,1)$.  Boundary sign-changing non-extreme points occur on the quarter circle $x^2+y^2=1$ with $x,y>0$, and at $(1,0)$ and $(0,1)$.
\end{answerbox}

\subsection*{Problem 4(iv): Global maximum and minimum on $[0,1]^2$}
From Problem 3(i), the global maximum is
\[
\frac18
\]
at
\[
\left(\frac12,\frac12\right).
\]
The global minimum is
\[
-1
\]
at
\[
(1,1).
\]

\begin{answerbox}{Final answer}
\[
\max_{[0,1]^2} f=\frac18
\quad\text{at}\quad \left(\frac12,\frac12\right),
\]
\[
\min_{[0,1]^2} f=-1
\quad\text{at}\quad (1,1).
\]
\end{answerbox}

\newpage
\section{Problem 5: Extreme values with constraints}

\begin{methodbox}{General method}
When the constraint can be solved directly, reduce the problem to one variable. When the constraint is geometric, use Lagrange multipliers:
\[
\grad f=\lambda \grad g.
\]
Then also check whether the constraint set is compact. If the constraint is not compact, a maximum or minimum may fail to exist.
\end{methodbox}



\begin{center}
\renewcommand{\arraystretch}{1.45}
\begin{tabular}{c}
\fbox{Constraint curve or surface: $g=0$}\\
At a constrained extremum, the level curve of $f$ is tangent to the constraint.\\
Therefore the normal directions agree: $\nabla f$ is parallel to $\nabla g$.\\
Hence the computational equation is $\boxed{\nabla f=\lambda\nabla g}$ together with $g=0$.
\end{tabular}
\end{center}



\subsection*{Problem 5(i)}
Find extreme values of
\[
f(x,y)=xy
\]
subject to
\[
x+y=1.
\]
Use the constraint to write
\[
y=1-x.
\]
Then
\[
f(x,1-x)=x(1-x)=x-x^2.
\]
Complete the square:
\[
x-x^2=-\left(x^2-x\right)
=-\left(\left(x-\frac12\right)^2-\frac14\right)
=\frac14-\left(x-\frac12\right)^2.
\]
Thus
\[
f(x,1-x)\le\frac14,
\]
with equality when
\[
x=\frac12.
\]
Then
\[
y=1-\frac12=\frac12.
\]
As $x\to\infty$ or $x\to-\infty$, the expression $x-x^2$ tends to $-\infty$. Hence there is no minimum on the whole line.

\begin{answerbox}{Answer}
Maximum value:
\[
\frac14\quad\text{at}\quad \left(\frac12,\frac12\right).
\]
There is no minimum, because the line $x+y=1$ is unbounded and $xy\to-\infty$ along it.
\end{answerbox}

\subsection*{Problem 5(ii)}
Find extreme values of
\[
f(x,y)=x^2+y^2
\]
subject to
\[
\frac{x}{a}+\frac{y}{b}=1.
\]
Assume $a\ne0$ and $b\ne0$.

Let
\[
g(x,y)=\frac{x}{a}+\frac{y}{b}-1.
\]
Then
\[
\grad f=(2x,2y),
\qquad
\grad g=\left(\frac1a,\frac1b\right).
\]
Lagrange multipliers give
\[
(2x,2y)=\lambda\left(\frac1a,\frac1b\right).
\]
So
\[
2x=\frac{\lambda}{a},
\qquad
2y=\frac{\lambda}{b}.
\]
Therefore
\[
x=\frac{\lambda}{2a},
\qquad
 y=\frac{\lambda}{2b}.
\]
Substitute into the constraint:
\[
\frac{1}{a}\cdot\frac{\lambda}{2a}+\frac{1}{b}\cdot\frac{\lambda}{2b}=1.
\]
Thus
\[
\frac{\lambda}{2a^2}+\frac{\lambda}{2b^2}=1.
\]
Factor out $\lambda/2$:
\[
\frac{\lambda}{2}\left(\frac1{a^2}+\frac1{b^2}\right)=1.
\]
Hence
\[
\lambda=\frac{2}{\frac1{a^2}+\frac1{b^2}}
=\frac{2a^2b^2}{a^2+b^2}.
\]
Therefore
\[
x=\frac{1}{2a}\cdot\frac{2a^2b^2}{a^2+b^2}
=\frac{ab^2}{a^2+b^2},
\]
\[
y=\frac{1}{2b}\cdot\frac{2a^2b^2}{a^2+b^2}
=\frac{a^2b}{a^2+b^2}.
\]
The minimum value is
\[
x^2+y^2
=\frac{a^2b^4}{(a^2+b^2)^2}+\frac{a^4b^2}{(a^2+b^2)^2}.
\]
Factor:
\[
x^2+y^2
=\frac{a^2b^2(a^2+b^2)}{(a^2+b^2)^2}
=\frac{a^2b^2}{a^2+b^2}.
\]
The constraint is a line, so it is unbounded. Along that line, $x^2+y^2\to\infty$. Hence there is no maximum.

\begin{answerbox}{Answer}
Minimum value:
\[
\frac{a^2b^2}{a^2+b^2}
\]
at
\[
\left(\frac{ab^2}{a^2+b^2},\frac{a^2b}{a^2+b^2}\right).
\]
There is no maximum.
\end{answerbox}

\subsection*{Problem 5(iii)}
Find extreme values of
\[
f(x,y)=x+y
\]
subject to
\[
\frac{x^2}{a^2}+\frac{y^2}{b^2}=1.
\]
Assume $a>0$ and $b>0$.

Let
\[
g(x,y)=\frac{x^2}{a^2}+\frac{y^2}{b^2}-1.
\]
Then
\[
\grad f=(1,1),
\qquad
\grad g=\left(\frac{2x}{a^2},\frac{2y}{b^2}\right).
\]
Lagrange multipliers give
\[
(1,1)=\lambda\left(\frac{2x}{a^2},\frac{2y}{b^2}\right).
\]
So
\[
1=\lambda\frac{2x}{a^2},
\qquad
1=\lambda\frac{2y}{b^2}.
\]
Therefore
\[
x=\frac{a^2}{2\lambda},
\qquad
 y=\frac{b^2}{2\lambda}.
\]
Substitute into the ellipse equation:
\[
\frac{1}{a^2}\left(\frac{a^2}{2\lambda}\right)^2+
\frac{1}{b^2}\left(\frac{b^2}{2\lambda}\right)^2=1.
\]
This becomes
\[
\frac{a^2}{4\lambda^2}+\frac{b^2}{4\lambda^2}=1.
\]
So
\[
\frac{a^2+b^2}{4\lambda^2}=1,
\]
which gives
\[
4\lambda^2=a^2+b^2.
\]
Thus
\[
2\lambda=\pm\sqrt{a^2+b^2}.
\]
If
\[
2\lambda=\sqrt{a^2+b^2},
\]
then
\[
x=\frac{a^2}{\sqrt{a^2+b^2}},
\qquad
 y=\frac{b^2}{\sqrt{a^2+b^2}}.
\]
The value is
\[
x+y=\frac{a^2+b^2}{\sqrt{a^2+b^2}}
=\sqrt{a^2+b^2}.
\]
If
\[
2\lambda=-\sqrt{a^2+b^2},
\]
then
\[
x=-\frac{a^2}{\sqrt{a^2+b^2}},
\qquad
 y=-\frac{b^2}{\sqrt{a^2+b^2}}.
\]
The value is
\[
x+y=-\sqrt{a^2+b^2}.
\]

\begin{answerbox}{Answer}
Maximum value:
\[
\sqrt{a^2+b^2}
\]
at
\[
\left(\frac{a^2}{\sqrt{a^2+b^2}},\frac{b^2}{\sqrt{a^2+b^2}}\right).
\]
Minimum value:
\[
-\sqrt{a^2+b^2}
\]
at
\[
\left(-\frac{a^2}{\sqrt{a^2+b^2}},-\frac{b^2}{\sqrt{a^2+b^2}}\right).
\]
\end{answerbox}

\newpage
\section{Problem 6: Lagrange multipliers}

\begin{methodbox}{General method}
For an equality constraint $g=0$, solve
\[
\grad f=\lambda \grad g,
\qquad g=0.
\]
After finding all candidates, evaluate $f$ at each candidate. If the constraint set is compact, the largest and smallest candidate values are the global maximum and minimum. If the constraint set is not compact, separately check whether a maximum or minimum can fail to exist.
\end{methodbox}

\subsection*{Problem 6(i)}
\[
f(x,y)=x^2+y^2,
\qquad g(x,y)=xy-1=0.
\]
The constraint is
\[
xy=1.
\]
Compute gradients:
\[
\grad f=(2x,2y),
\qquad
\grad g=(y,x).
\]
Lagrange equations are
\[
2x=\lambda y,
\qquad
2y=\lambda x,
\qquad
xy=1.
\]
Since $xy=1$, neither $x$ nor $y$ is zero. Divide the first equation by the second equation:
\[
\frac{2x}{2y}=\frac{\lambda y}{\lambda x}.
\]
So
\[
\frac{x}{y}=\frac{y}{x}.
\]
Multiplying by $xy$ gives
\[
x^2=y^2.
\]
Thus
\[
y=x
\quad\text{or}\quad y=-x.
\]
But $xy=1>0$, so $x$ and $y$ must have the same sign. Therefore
\[
y=x.
\]
Substitute into $xy=1$:
\[
x^2=1.
\]
Thus
\[
x=1\quad\text{or}\quad x=-1.
\]
So the candidates are
\[
(1,1),
\qquad (-1,-1).
\]
At both points,
\[
f=x^2+y^2=1+1=2.
\]
Also, by the AM-GM inequality,
\[
x^2+y^2\ge2|xy|=2
\]
on the constraint. Hence $2$ is the global minimum. There is no maximum because along $y=1/x$,
\[
f(x,1/x)=x^2+\frac1{x^2}\to\infty
\]
as $x\to\infty$.

\begin{answerbox}{Answer}
Minimum value $2$ at $(1,1)$ and $(-1,-1)$.  No maximum exists.
\end{answerbox}

\subsection*{Problem 6(ii)}
\[
f(x,y)=x^2+2y^2-4y,
\qquad g(x,y)=x^2+y^2-9=0.
\]
The constraint is the circle
\[
x^2+y^2=9.
\]
Compute gradients:
\[
\grad f=(2x,4y-4),
\qquad \grad g=(2x,2y).
\]
The Lagrange equations are
\[
2x=\lambda(2x),
\]
\[
4y-4=\lambda(2y),
\]
\[
x^2+y^2=9.
\]
The first equation gives
\[
(1-\lambda)x=0.
\]
So either $x=0$ or $\lambda=1$.

Case 1: $x=0$.  Then the constraint gives
\[
y^2=9,
\]
so
\[
y=3\quad\text{or}\quad y=-3.
\]
The points are $(0,3)$ and $(0,-3)$.

Case 2: $\lambda=1$.  The second equation becomes
\[
4y-4=2y.
\]
Thus
\[
2y=4,
\qquad y=2.
\]
The constraint gives
\[
x^2+2^2=9,
\qquad x^2=5.
\]
Thus
\[
x=\pm\sqrt5.
\]
The points are
\[
(\sqrt5,2),
\qquad (-\sqrt5,2).
\]

Evaluate $f$ at all candidates:
\[
f(0,3)=0+2(9)-12=6,
\]
\[
f(0,-3)=0+2(9)+12=30,
\]
\[
f(\pm\sqrt5,2)=5+2(4)-8=5.
\]
Since the circle is compact, the global minimum and maximum occur among these candidates.

\begin{answerbox}{Answer}
Minimum value $5$ at $(\sqrt5,2)$ and $(-\sqrt5,2)$.  Maximum value $30$ at $(0,-3)$.  The point $(0,3)$ is another constrained stationary point with value $6$.
\end{answerbox}

\subsection*{Problem 6(iii)}
\[
f(x,y)=y^2-x^2,
\qquad g(x,y)=\frac14x^2+y^2-1=0.
\]
The constraint is
\[
\frac{x^2}{4}+y^2=1.
\]
A direct parametrization is easiest:
\[
x=2\cos\theta,
\qquad y=\sin\theta.
\]
Then
\[
f(2\cos\theta,\sin\theta)=\sin^2\theta-(2\cos\theta)^2.
\]
So
\[
f=\sin^2\theta-4\cos^2\theta.
\]
Using $\sin^2\theta=1-\cos^2\theta$,
\[
f=1-\cos^2\theta-4\cos^2\theta=1-5\cos^2\theta.
\]
Since
\[
0\le\cos^2\theta\le1,
\]
we get
\[
-4\le f\le1.
\]
The maximum $1$ occurs when
\[
\cos^2\theta=0,
\]
so
\[
\cos\theta=0,
\qquad x=0,
\qquad y=\pm1.
\]
The minimum $-4$ occurs when
\[
\cos^2\theta=1,
\]
so
\[
y=0,
\qquad x=\pm2.
\]

\begin{answerbox}{Answer}
Maximum value $1$ at $(0,1)$ and $(0,-1)$.  Minimum value $-4$ at $(2,0)$ and $(-2,0)$.
\end{answerbox}

\subsection*{Problem 6(iv)}
\[
f(x,y)=xy,
\qquad g(x,y)=x^2+2y^2-6=0.
\]
Compute gradients:
\[
\grad f=(y,x),
\qquad \grad g=(2x,4y).
\]
Lagrange equations are
\[
y=2\lambda x,
\]
\[
x=4\lambda y,
\]
\[
x^2+2y^2=6.
\]
If $x=0$, then the second equation gives $0=4\lambda y$.  On the constraint, $y\ne0$, so this forces $\lambda=0$. But then the first equation gives $y=0$, contradiction. Hence $x\ne0$. Similarly $y\ne0$.

Substitute $y=2\lambda x$ into $x=4\lambda y$:
\[
x=4\lambda(2\lambda x)=8\lambda^2x.
\]
Since $x\ne0$,
\[
1=8\lambda^2.
\]
Thus
\[
\lambda=\pm\frac{1}{2\sqrt2}.
\]
From $y=2\lambda x$,
\[
y=\pm\frac{x}{\sqrt2}.
\]
Use the constraint:
\[
x^2+2\left(\frac{x^2}{2}\right)=6.
\]
Thus
\[
2x^2=6,
\qquad x^2=3.
\]
So
\[
x=\pm\sqrt3.
\]
Then
\[
y=\pm\sqrt{\frac32},
\]
with the sign depending on whether $y=x/\sqrt2$ or $y=-x/\sqrt2$.

If $x$ and $y$ have the same sign,
\[
xy=\sqrt3\sqrt{\frac32}=\frac3{\sqrt2}.
\]
If $x$ and $y$ have opposite signs,
\[
xy=-\frac3{\sqrt2}.
\]

\begin{answerbox}{Answer}
Maximum value:
\[
\frac3{\sqrt2}
\]
at
\[
\left(\sqrt3,\sqrt{\frac32}\right),
\qquad
\left(-\sqrt3,-\sqrt{\frac32}\right).
\]
Minimum value:
\[
-\frac3{\sqrt2}
\]
at
\[
\left(\sqrt3,-\sqrt{\frac32}\right),
\qquad
\left(-\sqrt3,\sqrt{\frac32}\right).
\]
\end{answerbox}

\subsection*{Problem 6(v)}
\[
f(x,y)=x^2y^2,
\qquad g(x,y)=x^2+y^2-1=0.
\]
Let
\[
u=x^2,
\qquad v=y^2.
\]
Then
\[
u\ge0,
\qquad v\ge0,
\qquad u+v=1.
\]
The objective becomes
\[
f=x^2y^2=uv.
\]
Since $v=1-u$,
\[
f=u(1-u)=\frac14-\left(u-\frac12\right)^2.
\]
Thus the maximum is
\[
\frac14
\]
when
\[
u=\frac12,
\qquad v=\frac12.
\]
This means
\[
x^2=y^2=\frac12,
\]
so
\[
x=\pm\frac1{\sqrt2},
\qquad y=\pm\frac1{\sqrt2}.
\]
All four sign choices are allowed.

The minimum occurs when
\[
uv=0.
\]
Thus either $u=0$ or $v=0$, meaning
\[
x=0\quad\text{or}\quad y=0.
\]
Together with $x^2+y^2=1$, the points are
\[
(\pm1,0),
\qquad (0,\pm1).
\]

\begin{answerbox}{Answer}
Maximum value $\frac14$ at
\[
\left(\pm\frac1{\sqrt2},\pm\frac1{\sqrt2}\right)
\]
with all four sign combinations. Minimum value $0$ at
\[
(\pm1,0),
\qquad (0,\pm1).
\]
\end{answerbox}

\subsection*{Problem 6(vi)}
\[
f(x,y)=\sqrt{xy},
\qquad g(x,y)=x+y-c=0.
\]
Assume $c>0$ and use the real-valued domain $xy\ge0$. On the line $x+y=c$, the relevant nonnegative segment is
\[
0\le x\le c,
\qquad y=c-x.
\]
Then
\[
f(x,c-x)=\sqrt{x(c-x)}.
\]
Since square root is increasing, maximizing $\sqrt{x(c-x)}$ is the same as maximizing
\[
x(c-x)=cx-x^2.
\]
Complete the square:
\[
cx-x^2=\frac{c^2}{4}-\left(x-\frac c2\right)^2.
\]
Therefore the maximum occurs at
\[
x=\frac c2.
\]
Then
\[
y=c-\frac c2=\frac c2.
\]
The maximum value is
\[
f\left(\frac c2,\frac c2\right)
=\sqrt{\frac c2\cdot\frac c2}
=\frac c2.
\]
The minimum on the closed nonnegative segment occurs at the endpoints:
\[
(x,y)=(0,c)
\quad\text{or}\quad (c,0).
\]
At both endpoints,
\[
f=0.
\]

\begin{answerbox}{Answer}
For $c>0$, maximum value $c/2$ at $(c/2,c/2)$; minimum value $0$ at $(0,c)$ and $(c,0)$.
\end{answerbox}

\subsection*{Problem 6(vii)}
\[
f(x,y)=x^2+y^2,
\qquad g(x,y)=x^4+y^4-1=0.
\]
Let
\[
u=x^2,
\qquad v=y^2.
\]
Then
\[
u\ge0,
\quad v\ge0,
\quad u^2+v^2=1.
\]
The objective is
\[
f=u+v.
\]
On the quarter circle $u^2+v^2=1$, $u,v\ge0$, the sum $u+v$ is minimized at the endpoints and maximized when $u=v$.

Minimum: endpoints
\[
(u,v)=(1,0)\quad\text{or}\quad (0,1).
\]
Then
\[
f=u+v=1.
\]
The corresponding points are
\[
(\pm1,0),
\qquad (0,\pm1).
\]

Maximum: $u=v$. Since
\[
u^2+v^2=1,
\qquad u=v,
\]
we get
\[
2u^2=1,
\qquad u=\frac1{\sqrt2}.
\]
Therefore
\[
f=u+v=\frac1{\sqrt2}+\frac1{\sqrt2}=\sqrt2.
\]
Since
\[
x^2=u=\frac1{\sqrt2},
\]
we get
\[
x=\pm 2^{-1/4}.
\]
Similarly
\[
y=\pm 2^{-1/4}.
\]

\begin{answerbox}{Answer}
Minimum value $1$ at $(\pm1,0)$ and $(0,\pm1)$.  Maximum value $\sqrt2$ at
\[
(\pm 2^{-1/4},\pm 2^{-1/4}),
\]
with all four sign choices.
\end{answerbox}

\subsection*{Problem 6(viii)}
\[
f(x,y)=x^4+y^4,
\qquad g(x,y)=x^2+y^2-1=0.
\]
Again let
\[
u=x^2,
\qquad v=y^2.
\]
Then
\[
u\ge0,
\quad v\ge0,
\quad u+v=1,
\]
and
\[
f=u^2+v^2.
\]
Substitute $v=1-u$:
\[
f=u^2+(1-u)^2.
\]
Expand:
\[
f=u^2+1-2u+u^2=2u^2-2u+1.
\]
Complete the square:
\[
2u^2-2u+1=2\left(u^2-u\right)+1
=2\left(\left(u-\frac12\right)^2-\frac14\right)+1.
\]
Thus
\[
f=2\left(u-\frac12\right)^2+\frac12.
\]
The minimum is
\[
\frac12
\]
when
\[
u=\frac12.
\]
Then
\[
v=\frac12.
\]
So
\[
x=\pm\frac1{\sqrt2},
\qquad y=\pm\frac1{\sqrt2}.
\]

The maximum occurs at endpoints $u=0$ or $u=1$. Thus
\[
f=1.
\]
The corresponding points are
\[
(\pm1,0),
\qquad (0,\pm1).
\]

\begin{answerbox}{Answer}
Minimum value $\frac12$ at
\[
\left(\pm\frac1{\sqrt2},\pm\frac1{\sqrt2}\right)
\]
with all four sign choices. Maximum value $1$ at $(\pm1,0)$ and $(0,\pm1)$.
\end{answerbox}

\subsection*{Problem 6(ix)}
\[
f(x_1,\ldots,x_n)=x_1+x_2+\cdots+x_n,
\]
subject to
\[
g(x_1,\ldots,x_n)=x_1^2+x_2^2+\cdots+x_n^2-1=0.
\]
This is optimization of a linear function on the unit sphere.

Write
\[
\ve{x}=(x_1,\ldots,x_n),
\qquad
\ve{1}=(1,\ldots,1).
\]
Then
\[
f(\ve{x})=\ip{\ve{x}}{\ve{1}}.
\]
The constraint is
\[
\norm{\ve{x}}=1.
\]
By the Cauchy-Schwarz inequality,
\[
|\ip{\ve{x}}{\ve{1}}|\le \norm{\ve{x}}\norm{\ve{1}}.
\]
Since
\[
\norm{\ve{x}}=1,
\qquad
\norm{\ve{1}}=\sqrt{n},
\]
we get
\[
|f(\ve{x})|\le\sqrt n.
\]
Therefore
\[
-\sqrt n\le f(\ve{x})\le\sqrt n.
\]
Equality in Cauchy-Schwarz occurs when $\ve{x}$ is a scalar multiple of $\ve{1}$.

For the maximum, take
\[
\ve{x}=\frac1{\sqrt n}(1,1,\ldots,1).
\]
Then
\[
f(\ve{x})=n\cdot\frac1{\sqrt n}=\sqrt n.
\]
For the minimum, take
\[
\ve{x}=-\frac1{\sqrt n}(1,1,\ldots,1).
\]
Then
\[
f(\ve{x})=-\sqrt n.
\]

\begin{answerbox}{Answer}
Maximum value:
\[
\sqrt n
\]
at
\[
\left(\frac1{\sqrt n},\ldots,\frac1{\sqrt n}\right).
\]
Minimum value:
\[
-\sqrt n
\]
at
\[
\left(-\frac1{\sqrt n},\ldots,-\frac1{\sqrt n}\right).
\]
\end{answerbox}

\subsection*{Problem 6(x)}
\[
H(p_1,\ldots,p_n)=\sum_{i=1}^n p_i\log_2 p_i,
\qquad p_i\ge0,
\]
subject to
\[
p_1+p_2+\cdots+p_n=1.
\]
We use the convention
\[
0\log_2 0=0,
\]
which is justified because
\[
\lim_{t\to0^+}t\log_2 t=0.
\]

For an interior critical point, assume $p_i>0$ for all $i$. Let
\[
g(p_1,\ldots,p_n)=p_1+\cdots+p_n-1.
\]
The derivative of $p_i\log_2p_i$ is
\[
\frac{d}{dp_i}\left(p_i\log_2p_i\right)
=\log_2p_i+\frac1{\ln2}.
\]
Lagrange multipliers give
\[
\log_2p_i+\frac1{\ln2}=\lambda
\]
for every $i$. Hence all $p_i$ are equal:
\[
p_1=p_2=\cdots=p_n.
\]
Since they sum to $1$,
\[
p_i=\frac1n
\]
for every $i$.

At the uniform point,
\[
H\left(\frac1n,\ldots,\frac1n\right)
=\sum_{i=1}^n \frac1n\log_2\left(\frac1n\right).
\]
There are $n$ identical terms, so
\[
H= n\cdot\frac1n\log_2\left(\frac1n\right)
=\log_2\left(\frac1n\right)
=-\log_2 n.
\]

Because the function $t\mapsto t\log_2t$ is convex on $[0,\infty)$, the sum $H$ is convex on the probability simplex. Therefore the interior critical point is the global minimum.

The maximum of a convex function on the simplex occurs on the boundary, and here it is reached at a vertex. If one probability equals $1$ and all others equal $0$, then
\[
H=1\cdot\log_2 1+0=0.
\]
Since $p_i\log_2p_i\le0$ for $0\le p_i\le1$, the value $0$ is the largest possible value.

\begin{answerbox}{Answer}
Minimum value:
\[
-\log_2 n
\]
at the uniform distribution
\[
p_1=p_2=\cdots=p_n=\frac1n.
\]
Maximum value:
\[
0
\]
at the vertices of the simplex, where one $p_j=1$ and all other $p_i=0$.
\end{answerbox}

\newpage
\section*{Final summary table}
\addcontentsline{toc}{section}{Final summary table}

\begin{center}
\renewcommand{\arraystretch}{1.35}
\begin{tabular}{|p{0.24\linewidth}|p{0.68\linewidth}|}
\hline
\textbf{Problem} & \textbf{Main result}\\
\hline
1 & Positive definite and negative definite criteria for $2\times2$ symmetric matrices; positive definite matrices have positive diagonal entries.\\
\hline
2 & Local extrema and saddle points obtained by solving $\grad f=0$ and applying the Hessian test.\\
\hline
3 & $xy(1-x^2-y^2)$ has global maximum $1/8$ at $(1/2,1/2)$ and global minimum $-1$ at $(1,1)$ on $[0,1]^2$; line-restricted minima need not imply a two-variable local minimum; least squares line is given by the normal equations.\\
\hline
4 & Detailed classification of $xy(1-x^2-y^2)$ on the open square and closed square.\\
\hline
5 & Constrained extrema solved by reducing constraints or using Lagrange multipliers.\\
\hline
6 & Lagrange multiplier candidates evaluated on each constraint; compactness determines whether global extrema must exist.\\
\hline
\end{tabular}
\end{center}

\end{document}
